\chapter*{Index}
\addcontentsline{toc}{chapter}{Index}
\newcommand{\ocindexpage}[1]{, p.~\csname pageref\endcsname{#1}}
\noindent Each name points at the statement, exercise or definition it names.

\begin{description}
\item[Bertrand's postulate] Theorem~2.5\ocindexpage{thm:bertrand}.
\item[best approximations by convergents] Lemma~9.4\ocindexpage{lem:best}.
\item[B\'ezout] Theorem~1.4\ocindexpage{thm:bezout}.
\item[Chebyshev-type bounds] Theorem~2.4\ocindexpage{thm:cheb}.
\item[Chinese remainder theorem] Theorem~3.5\ocindexpage{thm:crt}.
\item[continued-fraction error] Theorem~9.2\ocindexpage{thm:cfapprox}.
\item[convergent recurrence] Proposition~9.1\ocindexpage{prop:cfconv}.
\item[Dirichlet convolution] Proposition~4.3\ocindexpage{prop:conv}.
\item[Dirichlet's hyperbola formula] Theorem~4.6\ocindexpage{thm:hyperbola}.
\item[Eisenstein integers, the primes] Theorem~10.7\ocindexpage{thm:eisprimes}.
\item[Eisenstein's lemma] Lemma~6.6\ocindexpage{lem:eisenstein}.
\item[Euclid's lemma] Theorem~1.5\ocindexpage{thm:gauss}.
\item[Euclid--Euler theorem on even perfect numbers] Theorem~4.5\ocindexpage{thm:perfect}.
\item[Euclidean division] Theorem~1.1\ocindexpage{thm:div}.
\item[Euler's criterion] Theorem~6.2\ocindexpage{thm:eulercrit}.
\item[Euler's theorem, and Fermat's little theorem] Theorem~3.3\ocindexpage{thm:eulerthm}.
\item[Euler's bound on the prime reciprocals] Theorem~2.2\ocindexpage{thm:euler}.
\item[Fermat's descent, $x^4+y^4=z^2$] Theorem~8.2\ocindexpage{thm:fermat4}.
\item[Fermat's theorem for exponent $3$] Theorem~10.9\ocindexpage{thm:flt3}.
\item[fundamental theorem of arithmetic] Theorem~1.6\ocindexpage{thm:fta}.
\item[Gauss's circle estimate] Theorem~7.5\ocindexpage{thm:circle}.
\item[Gauss's lemma on the Legendre symbol] Theorem~6.3\ocindexpage{thm:gausslemma}.
\item[Gaussian integers, unique factorisation] Theorem~10.3\ocindexpage{thm:ufd}.
\item[Gaussian primes] Theorem~10.5\ocindexpage{thm:gaussprimes}.
\item[Hensel's lemma] Theorem~3.8\ocindexpage{thm:hensel}.
\item[infinitude of primes] Theorem~2.1\ocindexpage{thm:euclid-primes}.
\item[Jacobi symbol] Proposition~6.8\ocindexpage{prop:jacobi}.
\item[Jacobi's formula for $r_2(n)$] Theorem~10.6\ocindexpage{thm:r2}.
\item[Lagrange's four-square theorem] Theorem~7.4\ocindexpage{thm:foursq}.
\item[Lam\'e's bound] Theorem~1.8\ocindexpage{thm:lame}.
\item[law of quadratic reciprocity] Theorem~6.5\ocindexpage{thm:qr}.
\item[Legendre's formula for $v_p(n!)$] the count after the valuation rules\ocindexpage{prop:legendre}.
\item[Legendre's theorem on continued fractions] Theorem~9.5\ocindexpage{thm:legendre}.
\item[Markov's theorem] Theorem~8.6\ocindexpage{thm:markov}.
\item[M\"obius inversion] Theorem~4.4\ocindexpage{thm:mobius}.
\item[a Mordell equation, $y^2=x^3+7$] Theorem~8.5\ocindexpage{thm:mordell7}.
\item[order of a unit] Proposition~5.2\ocindexpage{prop:order}.
\item[Pell's equation] Theorem~9.8\ocindexpage{thm:pell}.
\item[periodic continued fractions] Theorem~9.6\ocindexpage{thm:periodic}.
\item[primitive Pythagorean triples] Theorem~8.1\ocindexpage{thm:pyth}.
\item[primitive roots modulo an odd prime] Theorem~5.3\ocindexpage{thm:cyclic}. The prime-power case is Theorem~5.5\ocindexpage{thm:cyclicn}.
\item[quadratic character of $2$] Corollary~6.4\ocindexpage{cor:two}.
\item[sums of two squares, a single prime] Theorem~7.2\ocindexpage{thm:twosq}. The general criterion is Theorem~7.3\ocindexpage{thm:twosqn}.
\item[Thue's lemma] Lemma~7.1\ocindexpage{lem:thue}.
\item[units modulo $2^k$] Theorem~5.6\ocindexpage{thm:twopower}. Squares modulo $8$ are Corollary~5.7\ocindexpage{cor:mod8squares}.
\item[Wilson's theorem] Theorem~3.4\ocindexpage{thm:wilson}.
\item[characters, the count] Theorem~11.4\ocindexpage{thm:t1-charcount}. Orthogonality is Theorem~11.5\ocindexpage{thm:t1-ortho}.
\item[Euler product for a character] Theorem~11.8\ocindexpage{thm:t1-eulerprod}.
\item[Legendre's sieve] Theorem~11.10\ocindexpage{thm:t1-sieve}.
\end{description}

\noindent The exercises below are the ones a later volume cites by number.
\begin{description}
\item[Legendre's formula, and the zeros of $1000!$] Exercise~1.3\ocindexpage{ex:zeros}.
\item[P\'olya's orchard] Exercise~1.15\ocindexpage{pb:orchard}.
\item[infinitely many primes $\equiv3\pmod4$] Exercise~2.1\ocindexpage{ex:mod4}.
\item[Furstenberg's topology] Exercise~2.8\ocindexpage{ex:furstenberg}.
\item[a prime divisor of $n!+1$] Exercise~2.9\ocindexpage{ex:t1-factorial-new-prime}.
\item[Lucas's theorem] Exercise~3.9\ocindexpage{ex:lucas}.
\item[$2^{\omega(n)}$] Exercise~4.2\ocindexpage{ex:muid}.
\item[primitive words] Exercise~4.7\ocindexpage{ex:primitive}.
\item[probability that two integers are coprime] Exercise~4.8\ocindexpage{ex:coprimeprob}.
\item[infinitely many primes $\equiv1\pmod n$] Exercise~5.7\ocindexpage{ex:cyclotomic}.
\item[the quadratic Gauss sum] Exercise~6.7\ocindexpage{ex:gausssum}.
\item[three squares modulo $8$] Exercise~7.1\ocindexpage{ex:three}.
\item[Liouville's number] Exercise~9.7\ocindexpage{ex:liouville}.
\item[{failure of unique factorisation in $\Z[\sqrt{-5}]$}] Exercise~10.2\ocindexpage{ex:zsqrt5}.
\end{description}

\noindent Symbols are listed by the place where the book defines them.
\begin{description}
\item[$a\mid b$] defined with the remainder conventions\ocindexpage{sym:divides}.
\item[$\gcd$] defined as the positive generator of $a\Z+b\Z$\ocindexpage{sym:gcd}.
\item[$\operatorname{lcm}$] defined with the valuations\ocindexpage{sym:lcm}.
\item[$v_p$] the valuation paragraph\ocindexpage{sym:valuation}.
\item[$\varphi$] the congruences chapter\ocindexpage{sym:phi}.
\item[$\mu$] the arithmetic-functions chapter\ocindexpage{sym:mu}.
\item[$\left(\frac{a}{p}\right)$] Definition~6.1\ocindexpage{def:legendre}.
\item[Jacobi symbol] Proposition~6.8\ocindexpage{prop:jacobi}.
\item[{$[a_0;a_1,a_2,\ldots]$}] Chapter~9\ocindexpage{ch:continued}.
\item[$\operatorname{ord}_m(a)$] Proposition~5.2\ocindexpage{prop:order}.
\item[$N(a+bi)$] Chapter~10\ocindexpage{ch:gaussian}.
\item[$\chi$, $L(s,\chi)$] Definition~11.1\ocindexpage{def:t1-character}.
\item[$|x|_p$] the field guide\ocindexpage{sym:padic}.
\end{description}
